\documentclass[unicode, 12pt, a4paper,oneside,fleqn]{article}

\usepackage{../src/styles} % Убедитесь, что путь к стилям верный
\usepackage{pgfplots}
\pgfplotsset{compat=1.18}
\usepackage{alltt}

\begin{document}

\begin{titlepage}
\begin{center}
\bfseries

{\Large Московский авиационный институт\\ (национальный исследовательский университет)
}

\vspace{48pt}

{\large Факультет информационных технологий и прикладной математики
}

\vspace{36pt}

{\large Кафедра вычислительной математики и~программирования
}

\vspace{48pt}

Лабораторная работа \textnumero 4 по курсу \enquote{Дискретный анализ}

\end{center}

\vspace{72pt}

\begin{flushright}
\begin{tabular}{rl}
Студент: & И.\,А. Казаков \\
Преподаватель: & А.\,А. Кухтичев \\
Группа: & М8О-206БВ-24 \\
Дата: & \\
Оценка: & \\
Подпись: & \\
\end{tabular}
\end{flushright}

\vfill

\begin{center}
\bfseries
Москва, \the\year
\end{center}
\end{titlepage}

\pagebreak

\CWHeader{Лабораторная работа \textnumero 4}

\CWProblem{
    Необходимо реализовать один из стандартных алгоритмов поиска образцов для указанного алфавита.
    
    \textbf{Вариант алгоритма:} Поиск одного образца, основанный на построении Z-блоков.
    
    \textbf{Вариант алфавита:} Слова не более 16 знаков латинского алфавита (регистронезависимые).
    
    Программа должна считывать искомый образец с первой строки, а затем обрабатывать текст, состоящий из слов и чисел, выводя координаты всех вхождений (номер строки и номер слова в строке).
}

\pagebreak

\section{Описание алгоритма}
Для решения задачи используется **Z-алгоритм**, адаптированный для работы с последовательностью слов (токенов). 

\subsection{Z-функция}
Z-функция от массива $S$ — это массив $Z$, где каждый элемент $Z[i]$ равен длине самого длинного общего префикса между массивом $S$ и его суффиксом, начинающимся с позиции $i$. 

Для поиска образца $P$ в тексте $T$ строится объединенный массив:
\[ S = P + [\$] + T \]
где $[\$]$ — специальный разделитель, не встречающийся в алфавите. Если для некоторого индекса $i$, относящегося к части текста $T$, значение $Z[i]$ равно длине образца $P$, то в этой позиции найдено вхождение.

\subsection{Особенности реализации}
\begin{itemize}
    \item \textbf{Токенизация:} Текст разбивается на отдельные слова с сохранением их координат (номер строки и порядковый номер слова в ней).
    \item \textbf{Регистронезависимость:} Все слова при сравнении приводятся к нижнему регистру.
    \item \textbf{Сложность:} Алгоритм работает за линейное время $O(N + M)$, где $N$ — количество слов в тексте, а $M$ — количество слов в образце.
\end{itemize}

\pagebreak

\section{Исходный код}
Реализация выполнена на языке C++.

\begin{lstlisting}[language=C++]
#include <algorithm>
#include <iostream>
#include <sstream>
#include <string>
#include <vector>

using namespace std;

struct WordInfo {
    int line;
    int pos;
};

void toLowerCase(string& s) {
    for (char& c : s) {
        if (c >= 'A' && c <= 'Z') {
            c = (char)(c - 'A' + 'a');
        }
    }
}

vector<int> computeZ(const vector<string>& s) {
    int n = s.size();
    vector<int> z(n, 0);
    int l = 0, r = 0;
    for (int i = 1; i < n; ++i) {
        if (i <= r) {
            z[i] = min(r - i + 1, z[i - l]);
        }
        while (i + z[i] < n && s[z[i]] == s[i + z[i]]) {
            z[i]++;
        }
        if (i + z[i] - 1 > r) {
            l = i;
            r = i + z[i] - 1;
        }
    }
    return z;
}

int main() {
    ios_base::sync_with_stdio(false);
    cin.tie(NULL);

    string patternLine;
    if (!getline(cin, patternLine)) {
        return 0;
    }

    vector<string> pattern;
    stringstream ssP(patternLine);
    string word;
    while (ssP >> word) {
        toLowerCase(word);
        pattern.push_back(word);
    }

    if (pattern.empty()) {
        return 0;
    }

    vector<string> textWords;
    vector<WordInfo> info;
    string textLine;
    int currentLine = 1;
    while (getline(cin, textLine)) {
        stringstream ssT(textLine);
        int currentWordPos = 1;
        while (ssT >> word) {
            string lowWord = word;
            toLowerCase(lowWord);
            textWords.push_back(lowWord);
            info.push_back({currentLine, currentWordPos});
            currentWordPos++;
        }
        currentLine++;
    }

    if (textWords.empty()) {
        return 0;
    }

    vector<string> combined;
    combined.reserve(pattern.size() + textWords.size() + 1);
    for (const auto& s : pattern) {
        combined.push_back(s);
    }
    combined.push_back("\0");
    for (const auto& s : textWords) {
        combined.push_back(s);
    }

    vector<int> z = computeZ(combined);
    int m = (int)pattern.size();

    for (int i = 0; i < (int)textWords.size(); ++i) {
        if (z[i + m + 1] == m) {
            cout << info[i].line << ", " << info[i].pos << "\n";
        }
    }

    return 0;
}

\end{lstlisting}

\pagebreak

\section{Исследование производительности}

Для анализа эффективности Z-алгоритма было проведено сравнение с наивным алгоритмом поиска (Naive Search). Исследование разделено на два сценария:

\begin{enumerate}
    \item \textbf{Normal Case:} Паттерн длиной 3 слова. Текст состоит из случайных строк. В этом случае наивный поиск работает быстро, так как несовпадение обнаруживается на первом же символе большинства слов.
    \item \textbf{Worst Case:} Паттерн длиной 100 слов (99 слов \enquote{a} и 1 слово \enquote{b}). Текст заполнен словами \enquote{a}. Здесь наивный поиск вынужден совершать $M-1$ сравнений для каждой позиции текста, что приводит к деградации сложности до $O(N \cdot M)$.
\end{enumerate}

\subsection{Результаты тестирования (Worst Case)}
Данные получены при фиксированной длине паттерна ($M=100$) и увеличении объема текста ($N$).

\begin{center}
\begin{tabular}{|c|c|c|}
\hline
Кол-во слов в тексте (тыс.) & Z-Algorithm (мс) & Naive Search (мс) \\ \hline
100 & 3.637 & 52.216 \\ \hline
200 & 6.954 & 102.203 \\ \hline
300 & 10.130 & 158.604 \\ \hline
400 & 13.350 & 202.200 \\ \hline
500 & 19.036 & 286.085 \\ \hline
\end{tabular}
\end{center}

\subsection{График зависимости времени от объема данных}

\begin{center}
\begin{tikzpicture}
    \begin{axis}[
        xlabel={Количество слов (тыс.)},
        ylabel={Время (мс)},
        grid=major,
        legend pos=north west,
        width=12cm, height=8cm,
        title={Worst Case: Z-Algorithm vs Naive Search}
    ]
    
    \addplot[color=blue, mark=circle] coordinates {
        (100, 3.637)
        (200, 6.954)
        (300, 10.130)
        (400, 13.350)
        (500, 19.036)
    };
    \addlegendentry{Z-Algorithm}

    \addplot[color=red, mark=x] coordinates {
        (100, 52.216)
        (200, 102.203)
        (300, 158.604)
        (400, 202.200)
        (500, 286.085)
    };
    \addlegendentry{Naive Search}
    \end{axis}
\end{tikzpicture}
\end{center}

\pagebreak

\section{Консоль}
Ниже представлен протокол сборки программы и запуска тестов производительности:

\begin{alltt}
# Сборка бенчмарка с оптимизацией O3
root@server:~/da_lab4# g++ -O3 benchmark.cpp -o solution

# Сценарий 1: Случайные данные (Normal Case)
root@server:~/da_lab4# python3 test_gen_normal.py 500000
root@server:~/da_lab4# ./solution < input.txt
Z-Algorithm:    14.999 ms (found 50)
Naive Search:   2.554 ms (found 50)

# Сценарий 2: Худший случай (Worst Case)
root@server:~/da_lab4# python3 test_gen.py 100000
Worst-case test generated: 100000 words, pattern length 100
root@server:~/da_lab4# ./solution < input.txt
Z-Algorithm:    3.637 ms (found 2)
Naive Search:   52.216 ms (found 2)

root@server:~/da_lab4# python3 test_gen.py 500000
Worst-case test generated: 500000 words, pattern length 100
root@server:~/da_lab4# ./solution < input.txt
Z-Algorithm:    19.036 ms (found 10)
Naive Search:   286.085 ms (found 10)
\end{alltt}

\section{Анализ результатов}
На основе полученных данных можно сделать следующие выводы:
\begin{enumerate}
    \item \textbf{Затраты на подготовку:} В обычном случае (случайный текст) Z-алгоритм проигрывает наивному поиску. Это связано с накладными расходами на выделение памяти под объединенный вектор \texttt{combined} и копирование строк.
    \item \textbf{Устойчивость к худшему случаю:} В сценарии с повторяющимися префиксами Z-алгоритм демонстрирует стабильную линейную сложность, в то время как время работы наивного алгоритма растет значительно быстрее из-за многократных вложенных сравнений.
    \item \textbf{Эффективность:} Z-алгоритм гарантирует время выполнения $O(N+M)$ вне зависимости от содержания текста, что делает его более надежным для обработки структурированных данных.
\end{enumerate}

\pagebreak

\section{Выводы}
В ходе выполнения лабораторной работы был реализован алгоритм поиска образца на основе построения Z-блоков. Алгоритм был успешно адаптирован для работы с последовательностями слов латинского алфавита. Тестирование подтвердило теоретическую оценку сложности $O(N + M)$ и показало преимущество метода над наивным поиском в условиях частых частичных совпадений образца с текстом.

\pagebreak

\input{../src/literature}

\end{document}